\section{Setup, configuration, and installation}

\subsection{Development prerequisites}
The production application requires Flutter, an Android SDK, and either an
Android emulator or a connected Android phone. Git and Java are provided by a
normal Flutter Android toolchain. The optional local sharing server additionally
requires a Docker-compatible runtime and the Supabase CLI. XeLaTeX and
\code{latexmk} reproduce this report.

From the repository root, verify the toolchain and dependencies:
\begin{verbatim}
C:\dev\flutter\bin\flutter.bat doctor -v
C:\dev\flutter\bin\flutter.bat pub get
C:\dev\flutter\bin\flutter.bat devices
\end{verbatim}

The device list must include an Android target. Start the complete application
with that exact identifier:
\begin{verbatim}
C:\dev\flutter\bin\flutter.bat run -d <android-device-id>
\end{verbatim}

No command must run before \code{main.dart} other than starting the Android
emulator or connecting the phone. In particular, a backend is not needed for
Explore, routes, Free Run, GPS tracking, quests, camera evidence, History,
achievements, Story Studio, PNG export, or the Android share sheet.

\subsection{Automated checks and Android build}
\begin{verbatim}
C:\dev\flutter\bin\dart.bat format --output=none --set-exit-if-changed lib test
C:\dev\flutter\bin\flutter.bat analyze
C:\dev\flutter\bin\flutter.bat test
C:\dev\flutter\bin\flutter.bat build apk --debug
\end{verbatim}

For the submission build, copy \path{android/key.properties.example} to the
ignored \path{android/key.properties}. Use the private keystore path and
credentials supplied to the team, then run:
\begin{verbatim}
C:\dev\flutter\bin\flutter.bat build apk --release
\end{verbatim}
The verified output is placed at \code{apk/app-release.apk}. Signing credentials
and the private keystore are never committed or packaged.

\subsection{Optional hosted sharing}
Only anonymous route-link sharing needs a server. After linking a Supabase
project, apply the migration, configure the function's server-side environment,
and deploy the function:
\begin{verbatim}
supabase login
supabase link --project-ref <project-ref>
supabase db push
supabase functions deploy share-run --no-verify-jwt
\end{verbatim}

Then pass only the public project URL and anonymous key to Flutter:
\begin{verbatim}
C:\dev\flutter\bin\flutter.bat run -d <android-device-id> `
  --dart-define=SUPABASE_URL=https://<project-ref>.supabase.co `
  --dart-define=SUPABASE_ANON_KEY=<public-anon-key>
\end{verbatim}

\subsection{Optional local sharing server}
For local function development, initialize once if the \code{supabase} directory
does not exist, start the local stack, apply the migration, and serve the
function:
\begin{verbatim}
supabase init
supabase start
supabase db reset --local
supabase functions serve share-run --no-verify-jwt
supabase status
\end{verbatim}

On the standard Android emulator, use
\code{http://10.0.2.2:54321} as \code{SUPABASE_URL}; a physical phone must use the
development computer's reachable LAN address. Cleartext HTTP is allowed only in
the debug Android manifest for this local workflow. Stop local services with
\code{supabase stop}. Hosted deployment remains the appropriate way to show a
link that another device can open.

\subsection{Focused browser harness}
Story Studio can be reviewed independently in a browser with:
\begin{verbatim}
C:\dev\flutter\bin\flutter.bat run -d edge `
  -t lib/features/story_studio/story_studio_demo.dart
\end{verbatim}
This is a development harness for the editor only. It does not represent the
complete production application or its SQLite/GPS/camera startup path.
